\section{Flow stages up close}

\begin{frame}[fragile]{Phase 1 -- Synthesis (Yosys)}
  \driving{Yosys}
  \begin{itemize}\small
    \item RTL $\rightarrow$ gate-level netlist using the GF180 standard-cell
          library (\cmd{gf180mcu\_fd\_sc\_mcu7t5v0}).
    \item Emits instance counts, estimated area, inferred FF / LUT shapes.
    \item Produces the first mapped \cmd{chip\_top.v} that feeds every
          downstream step.
  \end{itemize}
  \vspace{0.4em}
  Reference run numbers:
  \begin{lstlisting}[style=shellstyle]
design__instance__count,251615
design__instance__count__stdcell,84417
  \end{lstlisting}
  \note{%
    Yosys is fast (seconds) and deterministic. The 251\,k instance count
    looks huge -- it is dominated by $\sim$1100 padring cells and 81k tap
    cells inserted for well-tap density. The meaningful number is
    \cmd{instance\_count\_\_stdcell} (84k).\par
    If synthesis fails, it is almost always a Verilog issue
    (unresolved module, missing port, wrong bit width). LibreLane surfaces
    Yosys errors verbatim; grep the log for \cmd{ERROR:} to find the
    offending line.}
\end{frame}

\begin{frame}{Phase 2 -- Floorplan \& padring}
  \driving{OpenROAD}
  \begin{itemize}\small
    \item \cmd{OpenROAD.Floorplan} reads \cmd{DIE\_AREA} / \cmd{CORE\_AREA}
          and creates an empty chip frame.
    \item \cmd{OpenROAD.Padring} places I/O pads per \cmd{PAD\_SOUTH/...} order.
    \item \cmd{ODB.ManualMacroPlacement} drops the SRAMs at the coordinates
          from \cmd{config.yaml}.
    \item \cmd{CutRows} + \cmd{TapcellInsertion} prepare standard-cell rows.
  \end{itemize}
  \vspace{0.3em}
  \screenshotplaceholder{slides/figures/floorplan.png}{%
    OpenROAD: post-floorplan view (padring + SRAM placeholders).}
  \note{%
    This is the phase that turns YAML into geometry. If the pad list is
    wrong, the padring does not close and floorplan fails loudly.
    If \cmd{CORE\_AREA} is too tight, \cmd{CutRows} creates too few rows and
    placement fails later. If an SRAM instance coordinate overlaps the
    padring, \cmd{ManualMacroPlacement} fails here.\par
    Tap cells (well ties) are required by the PDK at a fixed pitch. Forget
    them and downstream LVS complains. LibreLane inserts them automatically.}
\end{frame}

\begin{frame}{Phase 3 -- Power Distribution Network}
  \driving{OpenROAD}
  \begin{itemize}\small
    \item OpenROAD runs \cmd{pdn\_cfg.tcl} to draw the grid.
    \item \term{Followpin} wires: thin VDD/VSS along every std-cell row.
    \item \term{Straps}: Metal2 vertical + Metal3 horizontal across the core.
    \item \term{Core ring}: loop around the core that connects to pads.
    \item Per-macro grids stitch SRAM power pins to the main grid.
    \item Checker target: \cmd{design\_\_power\_grid\_violation\_\_count = 0}.
  \end{itemize}
  \screenshotplaceholder{slides/figures/pdn.png}{%
    Core PDN with straps and ring visible.}
  \note{%
    Power is the one pass where you \emph{want} to touch Tcl. \cmd{pdn\_cfg.tcl}
    declares voltage domains and per-macro power grids; OpenROAD's PDN
    generator consumes it verbatim. Voltage domain naming must match
    \cmd{VDD\_NETS}/\cmd{GND\_NETS} in \cmd{config.yaml} or the tool refuses
    to start the phase.\par
    Signoff will re-check the grid (PDNSim + IR drop). A grid with zero
    \emph{construction} violations can still fail IR; if that happens,
    either widen straps or tighten pitch.}
\end{frame}

\begin{frame}[fragile]{Phase 4 -- Placement \& CTS}
  \driving{OpenROAD}
  \begin{itemize}\small
    \item \term{Global placement}: approximate positions (RePlAce).
    \item \term{Detailed placement}: legalise onto std-cell rows.
    \item \term{Clock Tree Synthesis}: buffer tree from \cmd{clk\_pad/Y}
          to every flop, minimising skew.
    \item \term{Post-CTS resizing}: fixes hold violations introduced by CTS.
  \end{itemize}
  \vspace{0.3em}
  Reference run:
  \begin{lstlisting}[style=shellstyle]
design__instance__count__class:clock_buffer,30
clock__skew__worst_setup,0.319   # ns
clock__skew__worst_hold, -0.328  # ns
  \end{lstlisting}
  \note{%
    Placement is iterative. Global placement runs on a quadratic wirelength
    model; detailed placement snaps to legal rows. CTS is where a
    well-behaved flow earns its keep: 30 clock buffers, skew below half a
    nanosecond, no PLL needed for 25\,MHz.\par
    Common failure: ``too many flops, not enough rows'' -- placement
    congestion exceeds a threshold. Fix by enlarging \cmd{CORE\_AREA} or
    adjusting \cmd{PL\_TARGET\_DENSITY\_PCT} (lower = more spread, less
    congestion, bigger die).}
\end{frame}

\begin{frame}[fragile]{Phase 5 -- Routing}
  \driving{OpenROAD}
  \begin{itemize}\small
    \item \term{Global routing}: assigns nets to routing regions.
    \item \term{Detailed routing}: puts wires on physical tracks.
    \item \term{Antenna repair}: inserts diodes on long metals that would
          damage gates during plasma etching.
    \item Iterates until zero DRC.
  \end{itemize}
  \vspace{0.3em}
  \begin{lstlisting}[style=shellstyle]
route__drc_errors__iter:0, 75
route__drc_errors__iter:5,  0      # clean after 6 iterations
route__wirelength__estimated, 153555  # um
antenna__violating__nets, 0
  \end{lstlisting}
  \note{%
    Routing is where designs blow up their wall-time budgets. OpenROAD's
    TritonRoute iterates: global route, detailed route, DRC, repair,
    repeat. Our run converges in 6 iterations; busy designs can take
    20+ on a congested die.\par
    \term{Antennas} are a fabrication concern, not a timing one: long
    metal wires accumulate plasma charge during etching, which can punch
    through a gate oxide. Diodes bleed the charge. LibreLane's repair pass
    inserts them automatically.\par
    If routing fails to converge, the fix is almost always the floorplan
    (more core area, lower density) or the congestion-vs-density knob.}
\end{frame}

\begin{frame}{Phase 6 -- Signoff (geometry + electrical)}
  \driving{Magic + KLayout + Netgen + OpenSTA}
  \textbf{Geometry (the slow part):}
  \begin{itemize}\footnotesize
    \item \term{Magic DRC} + \term{KLayout DRC} cross-check.
    \item \term{KLayout XOR} -- Magic GDS vs KLayout GDS (\emph{0 diff}).
    \item \term{Antenna / Density / Sealring} checks.
  \end{itemize}
  \textbf{Electrical:}
  \begin{itemize}\footnotesize
    \item \term{OpenRCX} extracts parasitics (3 corners).
    \item \term{OpenSTA} runs timing across 9 PVT + extraction corners.
    \item \term{PDNSim} does IR drop.
    \item \term{Netgen LVS} -- extracted netlist vs synthesized Verilog.
    \item Zero DRC / LVS / setup / hold / antenna violations = green chip.
  \end{itemize}
  \note{%
    Signoff is the "is this really manufacturable?" gate. Two independent
    DRC engines (Magic and KLayout) cross-check each other; an XOR confirms
    their geometries match. Any mismatch is a show-stopper.\par
    Electrical: RCX extracts RC parasitics from the final GDS. STA then
    re-checks timing with those real parasitics at every corner. IR drop
    checks the power grid under dynamic load. LVS is non-negotiable -- it
    is the single check that proves the layout implements the circuit you
    synthesized.\par
    Every check here writes one metric into \cmd{metrics.csv}. The next
    section is about reading that file.}
\end{frame}
